\chapter{The Unequal Starting Line}

Outside a convention center in Las Vegas, a real-estate operator makes a
distinction that deserves to be kept. He did not grow up with money, he says,
but he did grow up within reach of resources. There were people he could ask
for help. The decisive act was noticing them, accepting the humility of the
request, and asking.

He calls the difference resourcefulness. The word catches something that a
balance sheet misses. Advice, an introduction, an unused room, a borrowed tool,
an honest correction, a first customer, or an hour from someone who knows the
terrain can remain economically inert until a person recognizes it and acts.
A resource is not useful merely because it exists nearby.

But the claim becomes false if we stretch it one sentence too far. Not everyone
has an informed person willing to answer. Not everyone can risk disclosing a
need. The same request is heard differently depending on the speaker's age,
accent, race, gender, disability, reputation, legal status, and place in the
room. A person caring for a parent cannot spend an evening at the event where
the useful contact happens. A worker with no savings cannot treat an unpaid
month as an experiment. A licensed profession cannot be entered by confidence
alone. Resourcefulness makes available resources usable; it does not make
access equal.

That correction does not weaken agency. It makes agency specific. Instead of
telling a person to believe harder, we can ask what is actually within reach,
what could be built, what requires another person's consent, and what must
change before a move becomes safe. An honest account of wealth begins there.

\section{A starting position has more than one number}

Chapter 1 defined enough through capabilities. The same method helps us see a
starting point. Cash matters, but so do time, health, knowledge, credentials,
tools, language, legal permission, geography, relationships, reputation, and
the losses a household can survive. Obligations belong on the list as well. A
person supporting two children and an ill parent does not possess the same
experimental budget as a person with the same income and no dependents.

Call the set of actions that a person can take without violating a necessary
boundary the person's \emph{feasible set}. It is not a measure of worth. It is
a description of present options. A move belongs in that set only when its
requirements can be met: the money can be found, the time exists, the skill is
sufficient, the law permits it, a counterparty will participate, and failure
does not destroy the material floor.

The set changes. Practice adds skill. Saving adds runway. Reliable work adds
proof. A move can place a person near more customers. A friendship can reveal
an opening. Childcare can release time. Treatment can restore capacity. A
credential can open a gate. A good institution can convert a trustworthy plan
into credit. Policy can make an option possible for thousands of people at
once. Illness, eviction, discrimination, war, a platform ban, or a family
emergency can shrink the set just as quickly.

This language helps us avoid two symmetrical mistakes. The first says that
circumstance decides everything, so effort is theater. The second says that
effort decides everything, so an outcome reveals character. The interviews do
not support either view. They show people acting forcefully inside conditions
they did not choose, often with help they did not create, and sometimes with a
timing advantage they readily admit.

Consider three forms of capital that rarely appear together in a headline:

\begin{description}[leftmargin=2.2em,style=nextline]
\item[Capacity] Health, attention, emotional stability, time, and freedom from
  immediate danger determine whether a person can use anything else.
\item[Capability] Practiced skill, judgment, language, credentials, and tools
  determine which problems the person can solve reliably.
\item[Access] Geography, relationships, trust, legal status, institutions, and
  distribution determine whether that capability can reach a real opportunity.
\end{description}

Money can improve all three, but it is not interchangeable with them. A large
loan cannot restore exhausted judgment. A useful skill with no route to a
buyer earns nothing. An introduction without capability may open a door once
and close it permanently. Good strategy starts by locating the binding
constraint rather than repeating the resource that already happens to be
abundant.

\section{Luck leaves fingerprints}

One of the young entrepreneurs in the Las Vegas interviews says that much of
his first trading success arrived when markets were unusually easy. He names
the monetary conditions around the pandemic before describing what he did next.
That admission makes his account more, not less, instructive. Favorable prices
helped create the opening; delegation and a well-treated team later helped him
build beyond his own capacity. Timing explains part of the first result. It
does not explain every later decision.

Another entrepreneur, who grew up in Soviet social housing and entered mobile
phone retail after the Soviet Union collapsed, describes his own history as
catching a wave. Trade was opening, phones were spreading, oil prices were
rising, and he was present. An Atlanta developer makes the point even more
plainly: successful businesses often owe a great deal to timing, and some very
wealthy people have also been lucky.

Luck is not the opposite of work. It changes the return to work. The same
effort applied before a market exists, after it is crowded, or while an
institution is closed can lead to different outcomes. Action can increase
exposure to good fortune: more conversations create more chances to meet a
buyer, more experiments create more chances to discover demand, and financial
reserves create more chances to remain present when a cycle turns. Yet exposure
does not guarantee the event, and the ability to keep taking chances is itself
unequally distributed.

Tom Brady tells a story that holds both sides together. In college he
complained to the psychologist Greg Harden that other players were receiving
more opportunities. Harden told him to stop organizing his attention around
what others received and to make the most of the few practice opportunities he
did receive. Brady says he did. He also describes a father who believed in him,
mentors who taught him, teammates who performed with him, and specialists who
helped him maintain his body.

The useful lesson is local: give full attention to the opportunity actually in
front of you. It is not a proof that the allocation was fair, that three chances
are always enough, or that people who remain unseen failed to prepare. A
winner's account can reveal how one person used an opening. It cannot show us
all the equally serious people to whom the opening never came.

This is why a success story should be read with five questions:

\begin{enumerate}
\item What did the person choose and practice?
\item Which people or institutions made the choice usable?
\item Which prior resources absorbed the delay and the failures?
\item What timing, geography, health, or legal conditions enlarged the return?
\item Who might have made a similar effort and remained outside the frame?
\end{enumerate}

The questions do not subtract credit. They separate a mechanism that another
person can test from a circumstance that cannot simply be copied.

\section{What remains when the cash disappears}

In Denver, James Dumoulin asks several business owners what they would do if
their bank accounts reached zero. Their answers arrive quickly. Start from
skills and expertise. Package marketing or branding as a service. Ask business
owners where they are losing time or money. Sell on behalf of someone who
already has an offer. Call friends, bring people together, and listen for a
recurring problem. One speaker would seek financing.

The thought experiment is useful because it distinguishes money from the human
capital that helped produce it. Someone who has learned to diagnose a customer
problem, make a credible offer, deliver the work, collect payment, and earn a
referral is not truly returning to the beginning when a balance vanishes. The
skill, judgment, proof, and relationships remain.

But the phrase \emph{start from zero} conceals them. The speakers are imagining
zero cash while retaining healthy bodies, commercial knowledge, phones,
transport, reputations, contacts, legal standing, and the confidence produced
by previous success. A first-time worker with no proof is solving a different
problem. So is a person whose skill requires equipment, a license, medication,
childcare, or a market that has disappeared. Financing can bridge an operating
gap only when repayment has a credible source; otherwise it converts a lack of
cash into a fixed claim on uncertain income.

John Ruiz's early legal career makes the stack visible. He says his father came
from Cuba with five children and was at one point homeless. Ruiz became the
first in the family to attend college, then went to law school. He began his
firm with an \$800 credit card his father had given him and, unable to pay office
rent, exchanged hours of legal work for space. The compressed version is that
he began with almost nothing. The more useful version names what was present:
a father's help, education, a license, the ability to work, a landlord or
colleague willing to barter, and enough trust for the exchange to begin. Ruiz's
initiative joined those pieces. It did not conjure them from air.

Education occupies the same conditional place. A short Dallas conversation
contrasts a technology worker who now sees online learning everywhere with
another worker who entered technology because that was what he had been
educated to do, not because he had chosen the life deliberately. The second
speaker then overcorrects, arguing that knowing the right people matters more
than competence. Relationships do matter, especially at a gate. But an
introduction cannot substitute for clinical judgment, engineering safety,
legal authorization, or the repeated ability to do the work.

The practical question is not \emph{college or no college?} It is: what must
this path prove? Some paths require a degree or license. Some reward a compact
portfolio, an apprenticeship, a trade qualification, or a verified result.
Some educational programs deliver capability and access; others sell debt and
confidence without either. Choose the least costly route that supplies the
necessary skill, signal, permission, and network for the actual work.

Portable capital is therefore not skill in isolation. It is capability that
can still reach a real need under changed conditions. The next chapter will
build that idea carefully. Here it is enough to notice that the richest form
of recovery capital often sits in a bundle: skill, proof, trust, tools, health,
and a modest runway.

\section{A network is an institution in miniature}

The Denver answers about starting again turn repeatedly toward other people.
One speaker would gather friends and ask what problems they were facing.
Another credits relationships that hold him accountable to his stated values.
Others describe paid groups that exposed them to operators at a larger scale.
The benefit was not the room's prestige by itself. It was faster correction,
specific advice, higher standards, and evidence that a result they had treated
as impossible was being produced nearby.

Mike Collins gives a quieter example. He says a small community bank financed
his first truck because its people knew their customers. The bankers were not
merely a distant scoring system; they shared a community in which behavior and
reputation were visible. That local knowledge could turn a thin formal record
into enough trust for a first asset. It could also reproduce local prejudice or
exclude anyone who had not been allowed to become known. Relationship lending
is powerful precisely because human judgment enters the decision. Human
judgment can recognize context and carry bias at the same time.

This is social capital in its useful sense: not a pile of famous contacts, but
a history of reciprocal reliability. It grows when someone contributes before
making a large request, does small work well, tells the truth when a result is
bad, introduces people carefully, and remembers that access belongs partly to
the person who grants it. A network is not something one extracts. It is a set
of claims that people are willing to honor because prior conduct made the next
exchange reasonable.

The unequal part remains. Family and school can provide a network before a
person knows the word. Money buys travel, events, neighborhoods, memberships,
and time in which relationships can form. Migration can place a person near
customers and separate that person from every trusted helper at home. A Dallas
law-enforcement worker calls moving from Mississippi to Dallas his best
financial decision because he perceived more opportunity there. The move was
an act of agency. The ability to move, absorb housing costs, rebuild support,
and remain safe was part of the starting position.

For someone outside a powerful room, the first relationship move need not be a
paid mastermind or an attempt to meet a celebrity. It may be a trade group, a
library workshop, a former colleague, a customer interview, a local lender, a
professional association, a volunteer project, an open-source contribution, or
a careful note to a person whose work has genuinely been studied. The standard
is not status. It is whether the relationship improves evidence, capability,
accountability, or access without requiring the person to pretend to be richer
than they are.

\section{Agency needs a loss budget}

The last Las Vegas interview follows a woman who says she was nannying when
COVID removed her jobs. Bills remained and she was down to her last
twenty-dollar bill. She did not want her income to vanish whenever another household
no longer needed her hours. After trying several online business models, she
found short-term rentals, reported fast early growth, and later taught the
model to others.

Her account shows adaptive search: a shock made the old dependence visible;
she tested alternatives; one produced demand; she concentrated on it. It also
shows why advice cannot be detached from risk. She suggests business credit as
a route for someone without cash. Credit does not erase the starting gap. It
moves the gap into the future and adds a repayment obligation. Short-term
rentals depend on leases, zoning, platform rules, insurance, occupancy,
maintenance, neighbors, and local demand. Adding properties multiplies losses
as readily as gains when the unit economics are wrong.

The lesson is not to borrow into the business that happened to work for her.
It is to make the next experiment small enough to survive. Use present income
to buy a tool or a test when possible. Confirm permission before committing.
Measure one unit after every cost. Preserve the material floor. Do not make a
dependent household collateral for an unproven model.

The real-estate operator who praised resourcefulness supplies his own warning.
When asked what he sacrificed, he names health, friends, family, and sanity. He
says he is grateful and would not change the path. That is his judgment about
his life. It is not an invoice another reader is required to pay. A strategy
that wins money by consuming every other form of capital may cross the line
Chapter 1 asked us to draw.

Agency is strongest when it includes the power to stop, recover, and try again.
The relevant question is not merely \emph{Can I do this?} It is \emph{What will
still be intact if this attempt fails?}

\section{Make an opportunity inventory}

An opportunity inventory is more honest than either a gratitude list or a list
of grievances because it records both usable openings and binding constraints.
Write one page under four headings.

\begin{enumerate}
\item \textbf{Available now.} List the skills, time, tools, cash, relationships,
  permissions, and low-cost experiments that can be used within seven days.
  Include the material floor that must not be risked.
\item \textbf{Buildable.} List capabilities, proof, reserves, health, language,
  credentials, and relationships that could be developed over the next six to
  eighteen months. Name the smallest repeatable practice for each.
\item \textbf{Counterparty dependent.} List every option requiring a customer,
  employer, lender, partner, landlord, platform, regulator, or family member.
  State what that person needs to believe, what evidence would justify the
  belief, and what a respectful no would mean.
\item \textbf{Structurally blocked.} Name the legal, geographic, health, safety,
  care, discrimination, or institutional constraint that personal effort cannot
  remove quickly. Decide whether it calls for an ally, a move, professional
  help, collective action, policy, protection, or a different path.
\end{enumerate}

Then choose one action from the first heading and one investment from the
second. Do not pretend the fourth heading is the first. Do not use the fourth
to avoid the first. Record which assumptions the action will test, the largest
acceptable loss, and the date on which the evidence will be reviewed.

This practice changes the meaning of accountability. Accountability is not the
claim that every result was personally deserved. It is the refusal to surrender
the next available decision while remaining truthful about the ground on which
that decision is made.

Among all the resources a person can build, practiced capability travels best.
It can cross an employer, a city, and sometimes an industry. It is never wholly
independent of health, trust, tools, or demand, but it can enlarge a feasible set
without waiting for a spectacular opportunity. The next chapter begins with
that quieter form of capital: a skill good enough for someone else to pay for.
